\documentclass[11pt,a4paper]{article}
\usepackage[utf8]{inputenc}
\usepackage[french]{babel}
\usepackage[T1]{fontenc}
\usepackage{amsmath}
\usepackage{amsfonts}
\usepackage{amssymb}
\usepackage{graphicx}
\usepackage{enumitem}
\usepackage{tikz}
\usepackage{minted}
\newcommand{\mip}[1]{\mintinline{python}{#1}}
\newcommand{\mis}[1]{\mintinline{sql}{#1}}
\newcommand{\mib}[1]{\mintinline{bash}{#1}}

\usepackage{hyperref}
\usepackage{pstricks,pst-tree,pst-math,pst-eucl,pstricks-add}
\usepackage[margin=1cm,bottom=1.5cm]{geometry}
\begin{document}
\fbox{\textsc{25NSIJ1PO1}}


\bigskip

\textbf{Exercice 1} (6 points)

\medskip

\textit{Cet exercice porte sur les protocoles réseaux, l’algorithmique et la POO.}


\textbf{Partie A}

\begin{enumerate}
\item On complète la table de routage de $R_1$ avec RIP.

\begin{center}
\begin{tabular}{|c|c|c|}
\hline
\multicolumn{3}{|c|}{Table de routage de $R_1$}\\
\hline
destination&prochain saut&distance\\
\hline
$R_2$&$R_2$&0\\
\hline
$R_3$&$R_2$&1\\
\hline
$R_4$&$R_4$&0\\
\hline
$R_5$&$R_5$&0\\
\hline
$R_6$&$R_5$&1\\
\hline
\end{tabular}
\end{center}

\item La route suivie par un paquet du $LAN_1$ à Internet est $LAN_1$-$R_1$-$R_5$-$R_6$-Internet avec un routage RIP.

\item On complète la table de routage de $R_1$ avec OSPF.

\begin{center}
\begin{tabular}{|c|c|c|}
\hline
\multicolumn{3}{|c|}{Table de routage de $R_1$}\\
\hline
destination&prochain saut&distance\\
\hline
$R_2$&$R_2$&10\\
\hline
$R_3$&$R_2$&11\\
\hline
$R_4$&$R_2$&12\\
\hline
$R_5$&$R_2$&13\\
\hline
$R_6$&$R_2$&14\\
\hline
\end{tabular}
\end{center}


\item La route suivie par un paquet du $LAN_1$ à Internet est $LAN_1$-$R_1$-$R_2$-$R_3$-$R_4$-$R_5$-$R_6$-Internet avec un routage OSPF.


\item Si $R_2$ tombe en panne, la route suivie par un paquet du $LAN_1$ à Internet est $LAN_1$-$R_1$-$R_5$-$R_6$-Internet avec un routage OSPF, de distance 101.
\end{enumerate}

\bigskip

\textbf{Partie B}



\begin{enumerate}[resume]
\item On complète le tableau pour déterminer l'adresse réseau.

\begin{center}
\begin{tabular}{|c|c|c|c|c|}
\hline
\multicolumn{5}{|c|}{Calcul d'une adresse de réseau}\\
\hline
machine (binaire)&11000000&10101000&00000001&01100100\\
\hline 
masque (binaire)&11111111&11000000&00000000&00000000\\
\hline
réseau (binaire)&11000000&10000000&00000000&00000000\\
\hline
réseau (déc.pointée)&192&128&0&0\\
\hline
\end{tabular}
\end{center}

\item On complète le tableau pour déterminer l'adresse de broadcast.

\begin{center}
\begin{tabular}{|c|c|c|c|c|}
\hline
\multicolumn{5}{|c|}{Calcul d'une adresse de broadcast}\\
\hline
réseau (binaire)&11000000&10000000&00000000&00000000\\
\hline 
masque (binaire)&11111111&11000000&00000000&00000000\\
\hline
complément du masque (binaire)&00000000&00111111&11111111&11111111\\
\hline
broadcast (binaire) & 11000000&10111111&11111111&11111111\\
\hline
réseau (déc.pointée)&192&191&255&255\\
\hline
\end{tabular}
\end{center}

\item Avec l'adresse 172.16.1.100/16 et le masque 255.255.0.0, on a l'adresse 172.16.0.0 du $LAN_1$, l'adresse de broadcast 172.16.255.255, et il y a 65534 adresses disponibles, de 172.16.0.1 à 172.16.255.254.

\item On complète la méthode \mip{masquer}.


%    def masquer(self, masque: str)->str:
%        """
%        Détermine le préfixe masqué de l'adresse,
%        le masque (décimal pointé) étant passé en
%        paramètre.
%        >>> add = IPv4('192.168.1.100')
%        >>> add.masquer('255.192.0.0')
%        '192.128.0.0'
%        """
%        tmp = []
%        ip = self.octets()
%        crible = IPv4(masque).octets()

\begin{minted}[linenos,firstnumber=13]{python}
        for i in range(4):
            # Opération booléenne :
            tmp.append(ip[i] & crible[i])
        return ".".join([str(k) for k in tmp])
\end{minted}

\item On complète la méthode \mip{adresse_suivante}.


%    def adresse_suivante(self, adresse_max:str)->str:
%        """
%        Détermine l'adresse décimale pointée suivant
%        immédiatement l'adresse courante, sous réserve
%        d'existence d'une adresse disponible
%        >>> add = IPv4('192.168.1.100')
%        >>> add.adresse_suivante('192.168.1.254')
%        '192.168.1.101'
%        >>> add = IPv4('192.168.1.255')
%        >>> add.adresse_suivante('192.168.255.254')
%        '192.168.2.0'
%        """
%        assert self.adresse < adresse_max
%        liste_courante = self.octets()
%        liste_suivante = list()
%        retenue = 1
%        for index in range(4):
%            somme = liste_courante[3 - index] + retenue

\begin{minted}[linenos,firstnumber=19]{python}
            valeur, retenue = somme%256, somme//256
            liste_suivante = [str(valeur)] + liste.suivante
\end{minted}

%        return '.'.join(liste_suivante)



\end{enumerate}

\bigskip

\textbf{Exercice 2} (6 points)


\medskip

\textit{Cet exercice porte sur la programmation Python et la récursivité.}

\medskip

\begin{enumerate}
\item Avec la situation \tikz[baseline=-1ex]{\draw(0,0)rectangle(2.4,0.6);\draw(0.6,0)--(0.6,0.6);\draw(1.2,0)--(1.2,0.6);\draw(1.8,0)--(1.8,0.6);\draw[dotted](0.3,0.3)circle(1mm);\draw[fill=black](0.9,0.3)circle(1mm);\draw[fill=black](1.5,0.3)circle(1mm);\draw[fill=black](2.1,0.3)circle(1mm);
\draw(0.3,0)node[below]{\small 1};\draw(0.9,0)node[below]{\small 2};\draw(1.5,0)node[below]{\small 3};\draw(2.1,0)node[below]{\small 4};}
on peut jouer 1 et obtenir \tikz[baseline=-1ex]{\draw(0,0)rectangle(2.4,0.6);\draw(0.6,0)--(0.6,0.6);\draw(1.2,0)--(1.2,0.6);\draw(1.8,0)--(1.8,0.6);\draw[fill=black](0.3,0.3)circle(1mm);\draw[fill=black](0.9,0.3)circle(1mm);\draw[fill=black](1.5,0.3)circle(1mm);\draw[fill=black](2.1,0.3)circle(1mm);
\draw(0.3,0)node[below]{\small 1};\draw(0.9,0)node[below]{\small 2};\draw(1.5,0)node[below]{\small 3};\draw(2.1,0)node[below]{\small 4};}
ou jouer 3 et obtenir \tikz[baseline=-1ex]{\draw(0,0)rectangle(2.4,0.6);\draw(0.6,0)--(0.6,0.6);\draw(1.2,0)--(1.2,0.6);\draw(1.8,0)--(1.8,0.6);\draw[dotted](0.3,0.3)circle(1mm);\draw[fill=black](0.9,0.3)circle(1mm);\draw[dotted](1.5,0.3)circle(1mm);\draw[fill=black](2.1,0.3)circle(1mm);
\draw(0.3,0)node[below]{\small 1};\draw(0.9,0)node[below]{\small 2};\draw(1.5,0)node[below]{\small 3};\draw(2.1,0)node[below]{\small 4};}.
 
\item On écrit le code de la fonction \mip{initialiser}.

\begin{minted}{python}
def initialiser(n):
    return [False] * n
\end{minted}

\item On complète le code de la fonction \mip{victoire}.

\begin{minted}[linenos]{python}
def victoire(tab):
    for etat_case in tab:
        if etat_case == False:
            return False
    return True
\end{minted}


\item On écrit le code de la fonction \mip{indice_premiere_case_occupee}.

\begin{minted}[linenos]{python}
def indice_premiere_case_occupee(tab):
    for i in range(len(tab)):
        if etat_case:
            return i
    return None
\end{minted}

\item On écrit le code de la fonction \mip{coup_valide}.

\begin{minted}[linenos]{python}
def coup_valide(tab, case):
    return case==0 or case in range(len(tab)) and case == indice_premiere_case_occupee(tab)+1
\end{minted}


\item On écrit le code de la fonction \mip{changer_case}.

\begin{minted}[linenos]{python}
def changer_case(tab, case):
    if coup_valide(tab, case):
        tab[case] = not tab[case]
    return tab
\end{minted}


\item On complète le code de la fonction \mip{vider}.

\begin{minted}[linenos]{python}
def vider(n):
    if n == 1:
        print('Vider case 1')
    elif n > 1:
        vider(n-2)
        print('Vider case '+str(n))
        remplir(n-2)
        vider(n-1)
\end{minted}


\item On donne l'affichage produit par l'appel \mip{vider(3)}.


\begin{minted}[linenos]{python}
>>> vider(3)
Vider case 1
Vider case 3
Remplir case 1
Vider case 2
Vider case 1
\end{minted}


\item On écrit le code de la fonction \mip{remplir}.

\begin{minted}[linenos]{python}
def remplir(n):
    if n == 1:
        print('Remplir case 1')
    elif n > 1:
        remplir(n-1)
        vider(n-2)
        print('Remplir case '+str(n))
        remplir(n-2)
\end{minted}

\item Chacune des deux fonctions récursives \mip{vider} et \mip{remplir} s'appelle elle-même deux fois et l'autre une fois donc on a une complexité exponentielle, il est donc déraisonnable de vouloir traiter des baguenaudiers de grande taille.
 
\end{enumerate}

\bigskip


\textbf{Exercice 3} (8 points)

\medskip

\textit{Cet exercice porte sur la programmation en Python, les bases de données
relationnelles, le langage SQL, les systèmes d’exploitation et la sécurisation des
communications.}

\medskip

\textbf{Partie A}

\begin{enumerate}
\item L'appel \mip{gen_mdp(8, True, True, False)} convient.
\item On complète les lignes 8 à 10 de la fonction \mip{gen_mdp}.

\begin{minted}[linenos,firstnumber=8]{python}
    minuscules = [chr(i) for i in range(97,123)]
    majuscules = [chr(i) for i in range(65,91)]
    caracteres_speciaux = [chr(i) for i in range(33,48)] + [chr(i) for i in range(58,65)]
\end{minted}


\item On complète les lignes 13 et suivantes de la fonction \mip{gen_mdp}.

\begin{minted}[linenos,firstnumber=12]{python}
    if cont_min:
        jeu_caracteres = jeu_caracteres + minuscules
    if cont_maj:
        jeu_caracteres = jeu_caracteres + majuscules
    if cont_spe:
        jeu_caracteres = jeu_caracteres + caracteres_speciaux
\end{minted}


\item On complète la ligne 21 de la fonction \mip{gen_mdp}.

\begin{minted}[linenos,firstnumber=20]{python}
    for i in range(longueur):
        mot_de_passe = mot_de_passe + jeu_caracteres[randint(0, len(jeu_caracteres)] 
\end{minted}


\item La fonction \mip{gen_mdp} construit un mot de passe de longueur donnée en choisissant les caractères dans les catégories sélectionnées, mais ne garantit pas que des caractères de chaque catégorie seront utilisés.

\end{enumerate}

\bigskip

\textbf{Partie B}

%compte(mot_de_passe,utilisateur,renouvellement,id_site)
%site(id,nom_site,url)

\begin{enumerate}[resume]
\item L'attribut \texttt{mot\_de\_passe} de la relation \texttt{compte} sert de clé primaire et doit donc être unique, ce qui empêche d'utiliser plusieurs fois le même mot de passe.
\item La requête \texttt{SELECT url FROM site;} convient.
\item La requête \texttt{UPDATE compte SET mot\_de\_passe = 'yhTS?d@UTJe' WHERE mot\_de\_passe = '@rDfohpj!\&';} effectue la mise à jour.

\item La requête \texttt{SELECT id\_site FROM compte WHERE renouvellement < '2024-03-20';} convient.
\item Le format AAAA-MM-JJ permet de comparer des dates en utilisant l'ordre lexicographique usuel. 
\item La requête \texttt{SELECT utilisateur, mot\_de\_passe FROM compte JOIN site ON compte.id\_site = site.id }
 
\texttt{WHERE nom\_site = 'Votremailp' ORDER BY renouvellement;} convient.
\item Le fait de séparer les données en plusieurs tables facilite les mises à jour et permet
de limiter les incohérences ; p.ex on peut avoir plusieurs comptes sur un même site, et on pourra modifier l'url une seule fois dans la table \texttt{site} plutôt qu'une fois pour chacun des comptes.
\end{enumerate}

\bigskip

\textbf{Partie C}



\begin{enumerate}[resume]
\item L'appel \mip{chiffrement('/home/Documents/gestionnaire.db', '/home/Perso/secret.db', '/home/Perso/cle')} convient, avec des chemins absolus ; avec des chemins relatifs au répertoire courant, on peut utiliser.
l'appel \mip{chiffrement('gestionnaire.db', '../Perso/secret.db', '../Perso.cle')}.


\item On calcule \mip{0xA3 ^ 0x59} en passant par le binaire.


\begin{center}
\begin{tabular}{|c|c|}
\hline
hexa & binaire \\
\hline
A3& 1010.0011\\
\hline
59& 0101.1001\\
\hline
FA&1111.1010\\
\hline
\end{tabular}
\end{center}

\item Pour prouver l'identité \mip{(a XOR b) XOR b = a}, on présente les quatre cas dans un tableau.

\begin{center}
\begin{tabular}{|c|c|c|c|}
\hline
\mip{a}&\mip{b}&\mip{a XOR b}& \mip{(a XOR b) XOR b}\\
\hline
0&0&0&0\\
\hline
0&1&1&0\\
\hline
1&0&1&1\\
\hline
1&1&0&1\\
\hline
\end{tabular}
\end{center}

On constate que la première et la dernière colonne sont identiques, ce qui prouve la propriété.
\item Puisque la même clé permet de chiffrer et de déchiffrer le fichier, Alice utilise un chiffrement symétrique.

\item Les permissions sont \mib{rw} (lecture et écriture) pour Alice, \mib{r} (lecture seule) pour son groupe (le groupe eleves) et aussi \mib{r} pour les autres. Cela signifie que n'importe qui peut lire le fichier \mib{secret.db} ; Alice pourrait changer les permissions de ce fichier avec \mib{chmod 600 secret.db} ou bien  \mib{chmod go-r secret.db}.

Cependant sans accès au fichier \mib{cle}, l'attaque cryptographique est difficile puisque n'importe quel fichier de même taille que \mib{secret.db} peut être reconstruit en choisissant bien le fichier \mib{cle}.
\end{enumerate}

\bigskip

\textbf{Partie D}



\begin{enumerate}[resume]
\item La préconisation P1 est suivie, comme on l'a vu à la question 1.

Pour P2, on a vu que la fonction \mip{gen_mdp} de la partie A ne garantissait pas que chaque catégorie de caractère serait utilisée, on peut conseiller à Alice d'améliorer ce point.

La préconisation P3 sort du cadre de cet exercice ; on peut espérer qu'Alice, sensibilisée à la cybersécurité, évitera une telle erreur de débutant. Malgré tout, Alice stocke sa base de données en clair dans \mib{/home/Documents/gestionnaire.db}, et on a vu question 17 qu'elle pouvait être négligente sur les permissions des fichiers. Clairement on a un problème ici.


La préconisation P4 est manifestement suivie, mais utiliser une solution libre déjà existante serait sans doute préférable, avec moins de risque de failles de sécurité que la solution \og maison \fg{} d'Alice.




\end{enumerate}





\end{document}